% Source: https://artofproblemsolving.com/community/c3h1846778_two_random_proofs
% Stable identity: aops-topic-1846778-q02
\textbf{(a)} For $y=x^2-2x-2$: at $x=-2,1,2,3$ we get $y=4+4-2=6$, $1-2-2=-3$, $4-4-2=-2$ and $9-6-2=1$. For $6x=18+y-y^2$: at $y=6,-3,-2,1$ the right side is $18+6-36=-12$, $18-3-9=6$, $18-2-4=12$ and $18+1-1=18$, which are $6x$ for $x=-2,1,2,3$. So all four points lie on both curves. The first parabola has a vertical axis and the second a horizontal axis, so they are different parabolas through the same four points.

\begin{asy}
usepackage("amsmath");
size(7.5cm);
defaultpen(fontsize(10pt));
import graph;
real f(real x){return x^2-2*x-2;}
draw((-3.4,0)--(4.8,0),gray(0.7),Arrow(5));
draw((0,-5.3)--(0,7.4),gray(0.7),Arrow(5));
draw(graph(f,1-sqrt(10.2),1+sqrt(10.2)),blue+linewidth(0.9));
real g(real y){return (18+y-y^2)/6;}
path h=(g(-5),-5);
for(int i=1;i<=80;++i){real y=-5+11.6*i/80; h=h--(g(y),y);}
draw(h,red+linewidth(0.9));
pair[] P={(-2,6),(1,-3),(2,-2),(3,1)};
for(pair p:P) dot(p);
label("$(-2,6)$",(-2,6),NE,fontsize(9pt));label("$(1,-3)$",(1,-3),SE,fontsize(9pt));label("$(2,-2)$",(2,-2),SE,fontsize(9pt));label("$(3,1)$",(3,1),E,fontsize(9pt));
label("$y=x^2-2x-2$",(1+sqrt(7)+0.1,4),E,blue);
label("$6x=18+y-y^2$",(-2.1,-5),W,red);
label("\mbox{}",(min(currentpicture,user=true)+max(currentpicture,user=true))/2,black); // keeps the text after the figure black when LaTeX typesets the labels inline
\end{asy}

\textbf{(b)} For every $k>0$ consider the ellipse
\[
x^2+ky^2=1+k .
\]
At each of the four points, $x^2=y^2=1$, so the left side is $1+k$: all four points lie on every one of these ellipses. Different values of $k$ give different ellipses, since the $x$-intercepts $\pm\sqrt{1+k}$ are different. So infinitely many ellipses (in particular, infinitely many conics) pass through the four points. The figure shows $k=\dfrac{1}{3}$, $k=1$ (a circle) and $k=3$.

\begin{asy}
usepackage("amsmath");
size(5cm);
defaultpen(fontsize(10pt));
draw((-2.5,0)--(2.5,0),gray(0.7),Arrow(5));
draw((0,-2.4)--(0,2.5),gray(0.7),Arrow(5));
draw(scale(sqrt(4/3),2)*unitcircle,blue+linewidth(0.8));
draw(scale(sqrt(2))*unitcircle,red+linewidth(0.8));
draw(scale(2,sqrt(4/3))*unitcircle,heavygreen+linewidth(0.8));
pair[] P={(1,1),(-1,1),(-1,-1),(1,-1)};
for(pair p:P) dot(p);
label("$k=\dfrac{1}{3}$",(0.25,2),NE,blue);
label("$k=1$",(-1.0,-1.0)*sqrt(2)*0.71+(-0.05,-0.15),SW,red);
label("$k=3$",(2,0),SE,heavygreen);
label("\mbox{}",(min(currentpicture,user=true)+max(currentpicture,user=true))/2,black); // keeps the text after the figure black when LaTeX typesets the labels inline
\end{asy}

\textbf{Remark.} A general conic $Ax^2+Bxy+Cy^2+Dx+Ey+F=0$ has six coefficients, which matter only up to a common factor, so five points (no three on a line) are needed to pin it down. A parabola must also satisfy the extra condition $B^2=4AC$, so in general five points lie on no parabola at all, and part (a) shows that four points may lie on more than one.
